\section{Exact finite certificate and claim boundary}
\label{sec:tpc51-certificate}

The proofs in this paper are analytic.  The accompanying certificate is
a deterministic regression test for the finite algebra used in those
proofs; it is not numerical evidence for the prime number theorem and
does not test a twin-prime statistic.

The standard-library Python program
\begin{center}
\texttt{experiments/tpc51\_certificate.py}
\end{center}
uses exact rational arithmetic and performs \(18{,}498\) assertions in
six groups:
\begin{enumerate}[label=\textup{(C\arabic*)},leftmargin=3.2em]
 \item \(3{,}825\) explicit phase-kernel root identities;
 \item \(1{,}260\) cofactor null-vector checks for rectangular
       Vandermonde synthesis matrices;
 \item \(1{,}930\) boundary-zero and near-duplicate conditioning checks;
 \item \(1{,}008\) exact positive-definiteness and moment checks for
       finite sampled monomial Grams;
 \item \(10{,}471\) principal-minor checks for the sharp
       dimension-times-deletion perturbation bound on finite Walsh
       banks; and
 \item four exact endpoint fraction identities.
\end{enumerate}
The normalized source hash is
\begin{center}
{\small\texttt{6c8fac5abc221e00bd885ac59159b08d3a644a39426ff23910a883172d86c78c}},
\end{center}
the semantic payload hash is
\begin{center}
{\small\texttt{c64511069ce4e590eec2ac973663446bb3a8f6e16f93543daaf5e2334979a332}},
\end{center}
and the certificate-core hash is
\begin{center}
{\small\texttt{c0be9331cddaec485b3ec560e0368cee89c30f0df3bf42b384c3dd427c09cfe8}}.
\end{center}
The committed output is
\texttt{experiments/tpc51\_certificate\_output.txt}.  Re-running
\begin{center}
\texttt{python experiments/tpc51\_certificate.py
        --output experiments/tpc51\_certificate\_output.txt}
\end{center}
must reproduce it byte for byte when the normalized source hash is
unchanged.

\subsection{What is certified}

The finite checks catch sign errors in the differential phase kernel,
failure of the cofactor null-vector identity, accidental singularity of
the test Grams, an incorrect sparse-deletion constant, and arithmetic
mistakes in the endpoint exponents.  Exact rationals avoid tolerance
choices and platform-dependent linear algebra.

\subsection{What is not certified}

The certificate does not prove the analytic uniformity in
\cref{thm:tpc51-row-equidistribution}; that theorem is proved from the
fixed-modulus prime number theorem and the squarefree progression
count.  It does not enumerate the full physical coefficient system,
compute the nonlinear profile in \eqref{eq:tpc51-profile-envelope-reminder},
or test cancellation in a coherent prime sum.  Most importantly, its
finite Gram examples are consistency tests, not a data-driven
substitute for the limiting-independence hypothesis of
\cref{thm:tpc51-limit-gram}.

Thus the computational and theorem claims have the same boundary as
the rest of the paper: exact L0 algebra and the stated L1 row-sampling
interface, with no L2 fixed-shift prime-pair conclusion.
